% Experimental protocol for the BISON manuscript.
% Author: fangtoast
% Email: fangtoast@foxmail.com
% Copyright (c) 2026 fangtoast. All rights reserved.

\section{Experimental Protocol}

The evaluation follows the same evidence chain throughout.
\textbf{RQ1} tests full-band prediction against dense and local-convolution
baselines on unseen structures. \textbf{RQ2} evaluates a
coordinate-conditioned operator at withheld interpolation points and on an
unseen path segment. \textbf{RQ3} tests whether candidates ranked by BISON
retain their paired advantage after equal-budget MPB replay. Forward accuracy
alone is not used to infer inverse-design performance.

\subsection{Dataset and Physical Solver}

The MPB-generated dataset contains 10,000 binary,
\(\pgroup\)-symmetric unit cells on a \(16\times16\) grid. Each sample
contains 10 transverse-electric bands at 31 ordered wave-vector points. The
fixed split comprises 6,875 training, 2,500 validation, and 625 test
structures. All forward, query, and inverse evaluations reuse these structure
identities. The solver uses \(\epsilon_{\mathrm{air}}=1.0\),
\(\epsilon_{\mathrm{mat}}=8.9\), lattice constant \(a=1\), resolution 64,
10 bands, and eigensolver tolerance \(10^{-7}\).

The split unit is a complete structure and its band matrix. Wave-vector rows
from one structure are not divided between training and test sets. Query
partitions are constructed within the same structure identities. Interpolation
and held-segment tests therefore probe wave-vector generalization without
altering the structural test set.

\subsection{Forward and Query Comparisons}

The formal forward comparison includes BISON, a wide residual MLP, a lifted
CNN, and a plain MLP. Each model predicts the complete \(31\times10\) band
matrix from the input representation specified by the fair-comparison
protocol. Learning rates are selected on the validation split. Bands-2--10
test MRE is reported as the mean and standard deviation over seeds 0, 1, and
2. Each run also records parameter count, training time, inference latency,
and peak memory.

For the query task, each input pairs a structure with a physical wave vector,
and the output is a 10-band vector. We compare the FNO query model, DeepONet,
and an MLP under three protocols: the full path for unseen structures
(all-\(k\)), withheld interpolation points, and a contiguous held-out path
segment. The latter two protocols separate interpolation from path
extrapolation.

All model selection uses validation data, and the test set is reserved for the
reported estimates. The same bands-2--10 MRE definition is used for each
full-band comparison. Query evaluation reports band-order violations
separately because pointwise error and spectral ordering measure different
properties.

\subsection{Inverse Replay Protocol}

The inverse evaluation reuses frozen seed-0 BISON and query checkpoints; no
test target participates in checkpoint selection. All 625 test band matrices
receive a fixed pool of 128 candidates and a Top-5 MPB replay budget. The
single-mode routing-free generator is compared with the same decoder configured
with four cyclic mode embeddings. Both variants use latent width 32, hidden
width 256, depth four, identical objective weights, checkpoint-selection
rules, and replay budgets. Earlier retrieval and direct-inverse results are
excluded because the complete-test extension did not replay them.

This controlled comparison tests whether multiple fixed mode embeddings
improve the implemented latent decoder. It does not compare broader
generative-model families. The single-mode variant samples a continuous latent
distribution, and BISON provides the candidate ranking.

MPB replay uses eight isolated worker processes, failure-on-error semantics,
frequency-ascending band order, and structure-hash caching. Evaluation reports
the best of five replayed candidates per target. Confidence intervals use
10,000 bootstrap resamples, and the paired comparison includes a two-sided
sign test.

\subsection{Metrics and Reproducibility Boundary}

Forward results use the percentage MRE in Eq.~\eqref{eq:mre}. Inverse metrics
include surrogate MRE, real-MPB MRE, their difference in percentage points,
median replay MRE, and the effective unique candidate count. Runtime is
separated into candidate sampling, surrogate ranking, and solver replay.
Formal runs were executed in WSL Ubuntu with the
\texttt{meep\_env} Conda environment, PyTorch/CUDA in FP32 on an NVIDIA RTX
4060 Laptop GPU; MPB was invoked through the same environment. Appendix
\ref{app:provenance} records the artifact hashes and exact claim boundary.
